\documentclass[12pt,a4paper]{report}

% ---------- Encoding & language ----------
\usepackage[utf8]{inputenc}
\usepackage[T1]{fontenc}
\usepackage[english]{babel}

% ---------- Layout ----------
\usepackage[margin=1in]{geometry}
\usepackage{setspace}
\onehalfspacing

% ---------- Math ----------
\usepackage{amsmath,amssymb,amsthm}

% ---------- Graphics & tables ----------
\usepackage{graphicx}
\usepackage{booktabs}
\usepackage{float}

% ---------- Links & references ----------
\usepackage[hidelinks]{hyperref}
\usepackage{cleveref}

% ---------- Custom commands ----------
\newcommand{\R}{\mathbb{R}}

% ---------- Title-page fields ----------
\newcommand{\thesissupervisor}{Prof.\ Christopher Ehrmann}
\newcommand{\thesisschool}{School of Mechanical Engineering and Robotics}
\newcommand{\thesisprogram}{Sino-German Mechanical and Energy Engineering (CDHK)}

% ---------- Title metadata ----------
\title{Thermal Debonding of Adhesive Joints in EV Battery Pack Disassembly\\
\large A parametric simulation study of localized heating windows for releasing bonded cells without thermal runaway}
\author{Rizzolio Matteo}
\date{\today}

\begin{document}

% ============================================================
% FRONT MATTER
% ============================================================

\begin{titlepage}
\centering
\vspace*{0.4cm}

\includegraphics[width=7.2cm]{figures/tongji_logo.png}\par

\vspace{0.8cm}
{\LARGE\bfseries TONGJI UNIVERSITY}\par
\vspace{0.35cm}
{\large \thesisschool}\par
{\large \thesisprogram}\par

\vspace{2.0cm}
{\large\scshape Bachelor's Thesis}\par

\vspace{0.8cm}
{\LARGE\bfseries Thermal Debonding of Adhesive Joints\\[0.35cm]
in EV Battery Pack Disassembly}\par

\vspace{0.9cm}
{\large\itshape A parametric simulation study of localized heating windows\\
for releasing bonded cells without thermal runaway}\par

\vspace{1.2cm}
{\normalsize In partial fulfilment of the requirements for the degree of\\
Bachelor of Engineering}\par

\vspace{2.0cm}
\begin{tabular}{@{}r@{\hspace{0.6cm}}l@{}}
\textbf{Author:} & Matteo Rizzolio \\[0.35cm]
\textbf{Student ID:} & 2456368 \\[0.35cm]
\textbf{Supervisor:} & \thesissupervisor \\
\end{tabular}

\vfill
{\large Shanghai, China}\par
\vspace{0.2cm}
{\large October 2026}\par
\vspace*{0.6cm}
\end{titlepage}

\pagenumbering{roman}

\begin{abstract}
% Abstract text.
\end{abstract}

\tableofcontents
\listoffigures
\listoftables

\clearpage
\pagenumbering{arabic}

% ============================================================
% CHAPTER 1
% ============================================================

\chapter{Introduction}

\section{Background and motivation}
% Paragraph.

\section{Problem statement}
% Paragraph.

\section{Aim and scope of the thesis}
% Paragraph.

\section{Research questions}
% Paragraph.

\section{Structure of the thesis}
% Paragraph.

% ============================================================
% CHAPTER 2
% ============================================================

\chapter{Fundamentals and State of the Art}

\section{Adhesive bonding in battery packs}
% Paragraph.

\section{Thermal debonding methods}
% Paragraph.

\subsection{Localized heating methods: laser, induction and resistive heating}
% Paragraph.

\subsection{Alternative debonding approaches: ultrasonic, cryogenic and reversible adhesives}
% Paragraph.

\section{Thermomechanical behaviour of structural adhesives}
% Paragraph.

\section{Lithium-ion cell thermal safety}
% Paragraph.

\section{Battery dismantling, automation and regulation}
% Paragraph.

\section{Research gap and positioning of this work}
% Paragraph.

% ============================================================
% CHAPTER 3
% ============================================================

\chapter{Reference Case and Material Data}

\section{Reference case definition}
% Paragraph.

\section{Geometry and layer stacking}
% Paragraph.

\section{Material properties}
% Paragraph.

\subsection{Adhesive}
% Paragraph.

\subsection{Cell and aluminium adherends}
% Paragraph.

\subsection{Temperature-dependent properties}
% Paragraph.

\section{Operational temperature criteria}
% Paragraph.

\subsection{Adhesive debond threshold $T_{\mathrm{debond}}$}
% Paragraph.

\subsection{Cell safety criterion $T_{\mathrm{safe}}$}
% Paragraph.

\section{Summary of assumptions and data sources}
% Paragraph.

% ============================================================
% CHAPTER 4
% ============================================================

\chapter{Theoretical Modelling}

\section{Governing equations of transient heat conduction}
% Paragraph.

\section{Boundary conditions}
% Paragraph.

\subsection{Gaussian surface heat flux}
% Paragraph.

\subsection{Convection and radiation losses}
% Paragraph.

\section{Interface and contact resistance treatment}
% Paragraph.

\section{Dimensional analysis}
% Paragraph.

\subsection{Thermal penetration depth}
% Paragraph.

\subsection{Fourier number}
% Paragraph.

\subsection{Biot number}
% Paragraph.

\section{Definition of the process window}
% Paragraph.

% ============================================================
% CHAPTER 5
% ============================================================

\chapter{Finite-Element Model and Verification}

\section{Model geometry and domain extent}
% Paragraph.

\section{Symmetry assumptions and dimensionality}
% Paragraph.

\section{Mesh generation and interface refinement}
% Paragraph.

\section{Mesh-independence study}
% Paragraph.

\section{Energy-balance check}
% Paragraph.

\section{Verification against the semi-infinite-solid solution}
% Paragraph.

\section{Baseline case results}
% Paragraph.

% ============================================================
% CHAPTER 6
% ============================================================

\chapter{Parametric Study and Sensitivity Analysis}

\section{Parameter selection and ranges}
% Paragraph.

\section{Uncertain material and interface inputs}
% Paragraph.

\section{Fractional-factorial screening}
% Paragraph.

\section{Local refinement near the feasible-window boundary}
% Paragraph.

\section{Process-window map: heat flux versus duration}
% Paragraph.

\section{Sensitivity analysis}
% Paragraph.

\section{Summary of the parameter study}
% Paragraph.

% ============================================================
% CHAPTER 7
% ============================================================

\chapter{Discussion}

\section{Interpretation of the process window}
% Paragraph.

\section{Influence of assumptions and model limitations}
% Paragraph.

\section{Comparison with prior work}
% Paragraph.

\section{Implications for disassembly-oriented design}
% Paragraph.

\section{Threats to validity}
% Paragraph.

% ============================================================
% CHAPTER 8
% ============================================================

\chapter{Conclusions and Outlook}

\section{Main findings}
% Paragraph.

\section{Answers to the research questions}
% Paragraph.

\section{Limitations}
% Paragraph.

\section{Future work}
% Paragraph.

% ============================================================
% BACK MATTER
% ============================================================

\begin{thebibliography}{99}
% References.
\end{thebibliography}

\appendix

\chapter{Model Parameters}
% Paragraph.

\chapter{Supplementary Mesh-Study Data}
% Paragraph.

\end{document}
